%! Dividing Polynomials; Remainder and Factor Theorems — Unit 3, Lesson 3.3 — Exit Ticket (Answer Key)
\documentclass[10pt]{article}
\usepackage{algebra2-article}
\usepackage{algebra2-key}   % pulls in -boxes; do NOT also load -boxes
\newcommand{\TallMath}[1]{$\displaystyle #1\rule[-1.4em]{0pt}{3.2em}$}

\begin{document}

\pageheader{Unit 3, Lesson 3.3}{Exit Ticket --- Answer Key}
\namedateperiod

\vspace{0.10in}

No notes. Watch your placeholders and your signs.

\begin{enumerate}[leftmargin=1.8em, itemsep=12pt]
  \item \textbf{Synthetic division, then finish.} Divide $f(x) = x^3 + 2x^2 - 5x - 6$ by $x + 1$.

    \vspace{3pt}
    The divisor is $x+1$, so $r = $ \ans{$-1$}
    \begin{center}
    \renewcommand{\arraystretch}{1.6}
    \begin{tabular}{r|C{1.5cm}C{1.5cm}C{1.5cm}C{1.5cm}}
    \ans{$-1$} & \ans{$1$} & \ans{$2$} & \ans{$-5$} & \ans{$-6$} \\
        &     & \ans{$-1$} & \ans{$-1$} & \ans{$6$} \\ \cline{2-5}
        & \ans{$1$} & \ans{$1$} & \ans{$-6$} & \ans{$0$} \\
    \end{tabular}
    \end{center}

    \vspace{2pt}
    Quotient: \ans{$x^2+x-6$} \quad Remainder: \ans{$0$}

    \vspace{3pt}
    Factor the quotient and write $f$ completely factored: \ans{$(x+1)(x+3)(x-2)$}

    \vspace{3pt}
    The three zeros of $f$ are \ans{$-1,\ -3,\ 2$}

  \item \textbf{No dividing.} Let $g(x) = x^3 - 4x^2 + 2x + 5$.

    \vspace{3pt}
    What is the remainder when $g(x)$ is divided by $(x - 2)$? \ans{$1$}

    \vspace{3pt}
    Which theorem let you answer that without dividing? \ans{the Remainder Theorem}

    \vspace{3pt}
    Is $(x-2)$ a factor of $g$? \ans{No} \; How do you know? \ans{$g(2)=1\ne 0$, so the remainder is not zero.}

  \item \textbf{SOL-style multiple choice.} Which binomial is a factor of
    $p(x) = x^3 - 6x^2 + 11x - 6$?
    \begin{enumerate}[label=\Alph*., leftmargin=1.9em, itemsep=1pt]
      \item $x + 1$
      \item $x - 1$ \quad \textcolor{keyred}{\textbf{$\leftarrow$ correct}}
      \item $x - 6$
      \item $x + 6$
    \end{enumerate}
    \begin{itemize}[leftmargin=1.5em, nosep]
      \item By the Factor Theorem, test $r=1$: $p(1) = 1-6+11-6 = 0$, so $(x-1)$ is a factor.
      \item \textbf{A} is the sign-flip error --- $x+1$ means testing $r=-1$, and $p(-1) = -24$.
      \item \textbf{C} and \textbf{D} come from grabbing the constant term $6$ instead of testing anything: $p(6) = 60$ and $p(-6) = -504$.
    \end{itemize}
\end{enumerate}

\begin{teachernote}
Item~1 separates three different failures. A wrong $r$ (using $+1$) gives a bottom row ending in
$-8$, not $0$ --- that is the sign flip. A correct division that stops at $x^2+x-6$ is the
``forgot to finish'' pile, and it is the same habit Lesson~3.2 was built around. A correct
factorization with zeros listed as $1, 3, -2$ is a sign error in reading zeros off factors.

\vspace{3pt}
Item~2 is deliberately \emph{not} a factor. Students who divide it out anyway have the procedure but
not the theorem; students who answer $1$ in ten seconds have both. Item~3's distractor \textbf{A} is
the same sign flip as item~1, so a student who misses both has one problem, not two.

\vspace{3pt}
Sort into four piles --- \emph{got it} / \emph{sign flip on $r$} / \emph{placeholder or arithmetic}
/ \emph{stopped before factoring}. A large sign-flip pile is a two-minute fix at the top of
Lesson~3.4 (write $x+1 = x-(-1)$ every time). A large ``stopped early'' pile matters more, because
Lesson~3.4's whole method is test $\rightarrow$ divide $\rightarrow$ \emph{finish}.
\end{teachernote}

\end{document}
